% thesis_vi.tex — Luận văn tiếng Việt (bản SIC CRISP-DM)
% C2 Evasion RL: học tăng cường mức gói tin để né phát hiện C2.
% Biên dịch độc lập từ thư mục docs/:
%   pdflatex thesis_vi && bibtex thesis_vi && pdflatex thesis_vi && pdflatex thesis_vi
% (hoặc: make -C thesis ... nếu dùng Makefile sẵn có)
\documentclass[12pt,a4paper]{report}

% ---------------------------------------------------------------------------
% Preamble nội tuyến (rút gọn từ docs/thesis/preamble.tex để file tự chứa)
% ---------------------------------------------------------------------------
\usepackage[utf8]{inputenc}
\usepackage[T5]{fontenc}   % bảng mã T5 (vntex) — bắt buộc cho tiếng Việt
\usepackage{lmodern}
% Tiếng Việt phải là ngôn ngữ CHÍNH (đặt cuối) để kích hoạt T5;
% nếu không, lệnh \ohorn ... vẫn gắn OT1 và build sẽ lỗi.
\usepackage[english,vietnamese]{babel}
\usepackage{amsmath, amssymb, amsfonts}
\usepackage{graphicx}
\usepackage{hyperref}
\usepackage{listings}
\usepackage{xcolor}
\usepackage{fancyhdr}
\usepackage[margin=1in]{geometry}
\usepackage{natbib}
\usepackage{booktabs}
\usepackage{caption}
\usepackage{subcaption}
\usepackage{enumitem}
\usepackage{microtype}

\DeclareUnicodeCharacter{2713}{\ensuremath{\checkmark}}
\DeclareUnicodeCharacter{2717}{\ensuremath{\times}}

\newcommand{\eg}{\emph{e.g.}}
\newcommand{\ie}{\emph{i.e.}}
\newcommand{\etal}{\emph{et al.}}
\newcommand{\PPO}{\texttt{PPO}}
\newcommand{\IDS}{\texttt{IDS}}
\newcommand{\XGB}{\texttt{XGBoost}}
\newcommand{\RL}{\textit{RL}}
\newcommand{\MDP}{\textit{MDP}}
\newcommand{\CTU}{\textit{CTU-13}}
\newcommand{\Snort}{\textit{Snort}}

\lstset{
  basicstyle=\ttfamily\small, breaklines=true, columns=fullflexible,
  frame=single, numbers=left, numberstyle=\tiny, keywordstyle=\color{blue},
  commentstyle=\color{gray}, stringstyle=\color{red},
  backgroundcolor=\color{lightgray!10}, language=Python
}

\hypersetup{
  colorlinks=true, linkcolor=blue, citecolor=blue, urlcolor=blue,
  pdftitle={C2 Evasion RL: Học tăng cường mức gói tin cho né tránh phát hiện C2},
  pdfauthor={Le Hoang Phat}
}

\pagestyle{fancy}
\fancyhf{}
\rhead{\thepage}
\lhead{C2 Evasion RL}
\renewcommand{\headrulewidth}{0.4pt}
\setlength{\headheight}{15pt}

\selectlanguage{vietnamese}

\title{\textbf{C2 Evasion RL}\\[4pt]
Học tăng cường mức gói tin để né phát hiện lưu lượng C2 botnet\\
trên Snort 2.9.20 với bộ luật ET Open}
\author{Lê Hoàng Phát}
\date{\today}

\begin{document}

\maketitle

% ===========================================================================
\begin{abstract}
\noindent
Luận văn này nghiên cứu câu hỏi: \emph{một agent học tăng cường (RL) có thể tự động
biến đổi các gói tin beacon C2 thật để né một IDS chữ ký thật, trong khi luồng C2
vẫn còn hoạt động được không?} Chúng tôi xây dựng môi trường \texttt{HiddenDefenderEnv}
với 14 cơ chế biến đổi gói tin, dùng \texttt{MaskablePPO} để chọn cơ chế, và chấm điểm
mọi hành động bằng chính nhị phân \textbf{Snort 2.9.20} (bộ luật ET Open C2, 21\,374
luật) chứ không dùng mô hình xấp xỉ khi đo kết quả chính.

Kết quả chính được đo bằng Snort thật: trên tập test held-out 40 luồng, agent đạt
\textbf{40/40 né hợp lệ (100\%)} — không báo động, và ngữ nghĩa C2 vẫn nguyên vẹn —
trong khi \textbf{control âm tính} (luồng gốc chưa biến đổi) bị Snort báo động
\textbf{40/40 (tỷ lệ né 0\%)}. Đây là bằng chứng rằng né tránh là thật, không phải
lỗi đo lường.

Tuy nhiên, luận văn cũng báo cáo trung thực hai phát hiện \emph{phủ định} quan trọng:
(1) với defender tĩnh này, bài toán tìm kiếm gần như suy biến — mỗi framing chỉ có một
hành động tối ưu gần như cố định (\texttt{split8} cho luồng nhị phân, \texttt{split24}
cho luồng HTTP), nên RL \textbf{không vượt} được các baseline đơn giản; và (2) khi
kiểm tra chéo capture (cross-capture), tỷ lệ né trung bình chỉ \textbf{38{,}89\%}
(độ lệch chuẩn 43{,}74\%), và trên cả ba capture agent \textbf{hoà} với control
``corrupt-everything'' — nghĩa là nó không học được gì mà hành động tầm thường
``làm hỏng mọi gói'' không làm được. Một capture (win13) đạt 0\% vì các luồng của nó
gần như không có payload để biến đổi — đây là giới hạn của \emph{không gian hành động},
không phải kết luận về họ botnet.

Luận văn đóng gói kèm một quy trình đo lường nghiêm ngặt (một cặp (flow, cơ chế) cho
mỗi pcap, control âm tính bắt buộc, trần oracle, bộ verifier 81 kiểm tra tự động) và
ghi lại sáu lỗi phương pháp đã từng tạo ra kết quả sai trông rất hợp lý. Đóng góp của
công trình là sự trung thực của phép đo và việc chỉ ra rằng, dưới một defender tĩnh
đơn luật, RL chưa tạo thêm giá trị so với tìm kiếm tham lam.
\end{abstract}

\tableofcontents
\listoftables
\listoffigures

% ===========================================================================
\chapter{Giới thiệu \& Động lực}
\label{ch:intro}

\section{Bối cảnh}

Botnet là mạng lưới các máy bị nhiễm (bot) nhận lệnh từ một máy chủ điều khiển
(command-and-control, C2). Kênh C2 là điểm sống còn: cắt được nó là vô hiệu hoá cả
botnet. Vì vậy các hệ thống phát hiện xâm nhập (IDS) như \Snort{} \citep{snort_project}
với bộ luật chữ ký mở như Emerging Threats (ET) Open \citep{emerging_threats} tập
trung vào việc khớp các mẫu byte đặc trưng trong lưu lượng C2.

Đây là một trò chơi ``mèo vờn chuột'': phía tấn công liên tục tìm cách che giấu kênh
C2 (mã hoá, phân mảnh, chèn nhiễu), phía phòng thủ cập nhật luật để bắt lại. Việc tạo
ra các biến đổi né tránh một cách thủ công tốn nhiều công sức và đòi hỏi hiểu biết sâu
về giao thức. Câu hỏi đặt ra: \emph{có thể tự động hoá quá trình này bằng học tăng
cường không?}

\section{Phát biểu bài toán}

Cho một tập luồng C2 đã được gán nhãn (lấy từ các capture malware thật), mục tiêu là
huấn luyện một \emph{agent} sao cho, với mỗi luồng, nó chọn một phép biến đổi gói tin
để:

\begin{enumerate}[label=(\arabic*)]
  \item Snort \textbf{không} còn báo động trên luồng đó (\emph{né tránh}), và
  \item kênh C2 vẫn còn hoạt động được — lệnh C2 vẫn tới được endpoint
        (\emph{giữ ngữ nghĩa}).
\end{enumerate}

Hai điều kiện này phải cùng đúng; một phép biến đổi xoá luôn kênh C2 thì không được
tính là né tránh. Ngoài ra agent phải \emph{tối thiểu hoá chi phí} biến đổi, vì sửa
nhiều gói làm tăng nguy cơ bị phát hiện bằng kênh phụ và làm hỏng độ tin cậy của C2.

\section{Động lực hai chiều}

\begin{itemize}
  \item \textbf{Phía phòng thủ:} hiểu được cách một kẻ tấn công \emph{thực sự} có thể
        tạo lưu lượng trốn tránh giúp xây dựng luật phát hiện bền vững hơn.
  \item \textbf{Phía học thuật:} khoảng cách giữa ``chấm điểm bằng mô hình xấp xỉ''
        và ``chấm điểm bằng IDS thật'' là một hạn chế cơ bản của transfer learning
        trong bối cảnh đối kháng; đo lường nó một cách trung thực có giá trị riêng.
\end{itemize}

\section{Đóng góp}

\begin{enumerate}[label=\textbf{C\arabic*.}]
  \item \textbf{Đường ống đầu-cuối trên IDS thật.} Một vòng lặp huấn luyện RL trong đó
        mọi hành động được chấm bởi nhị phân Snort 2.9.20 thật (chế độ \texttt{--real-snort}),
        tích hợp với pcap thật từ Stratosphere IPS \citep{stratosphere_ctu13}.
  \item \textbf{Kết quả chính có kiểm chứng.} 40/40 luồng held-out né hợp lệ dưới Snort
        thật, đối chiếu control âm tính 40/40 bị báo động.
  \item \textbf{Kết quả phủ định được báo cáo thẳng thắn.} RL không vượt baseline tham
        lam, và kiểm tra chéo capture cho thấy policy \emph{hoà} với control
        ``corrupt-everything'' trên cả ba capture — nên con số né chéo capture
        \textbf{không} phải thành tựu của RL.
  \item \textbf{Kỷ luật đo lường.} Quy trình một cặp (flow, cơ chế) cho mỗi pcap,
        control âm tính bắt buộc, trần oracle, và bộ verifier tự động 81 kiểm tra
        (\texttt{controls/verify\_report\_claims.py}) đảm bảo mọi con số trong báo cáo
        đều truy được về file JSON gốc.
  \item \textbf{Bài học về lỗi phương pháp.} Sáu lỗi (trích xuất luồng một chiều, viết
        lại gói tin ngây thơ, kiểm tra ngữ nghĩa lỏng, chấm điểm theo lô, luật bị
        throttle, cửa sổ đọc alert quá ngắn) từng tạo ra kết quả sai trông rất hợp lý,
        và cách chúng bị phát hiện.
\end{enumerate}

\section{Phạm vi}

Kết quả được xác thực trên \textbf{bộ luật tĩnh ET Open}, trên \textbf{cùng một họ
malware} (một nhị phân, nhiều ngày). Luận văn \emph{không} tuyên bố tổng quát hoá
sang mọi họ botnet, cũng không tuyên bố chống lại IDS thích ứng. Những giới hạn này
được nêu rõ ở Chương~\ref{ch:discussion}.

% ===========================================================================
\chapter{Nghiên cứu liên quan}
\label{ch:related}

\section{Ví dụ đối kháng trong học máy}

Học máy đối kháng nghiên cứu cách tạo đầu vào khiến mô hình phân loại sai trong khi
vẫn giữ tính hợp lệ. \citet{biggio2013evasion} đặt nền móng cho \emph{tấn công né
tránh} (evasion attack); \citet{goodfellow2014explaining} và
\citet{carlini2016towards} phát triển các phương pháp gradient để tối ưu nhiễu đối
kháng. \citet{papernot2016limits} mở rộng sang miền rời rạc (văn bản, tệp nhị phân).
Điểm chung: né tránh là một \emph{bài toán tối ưu có ràng buộc}, và trong không gian
rời rạc (như gói tin) thì tìm kiếm trở nên khó hơn nhiều so với tối ưu liên tục.

Trong miền malware, \citet{pierazzi2019intriguing} và \citet{demetrio2020intriguing}
nhấn mạnh khái niệm \emph{problem-space} (không gian bài toán): một biến đổi phải vừa
qua được bộ phân loại, vừa giữ được chức năng thực thi của mẫu. Đây chính xác là ràng
buộc ``giữ ngữ nghĩa C2'' mà chúng tôi áp dụng.

\section{Phát hiện C2 và IDS chữ ký}

\Snort{} \citep{snort_project} là IDS dựa trên chữ ký phổ biến, khớp các mẫu byte và
điều kiện header. Bộ luật ET Open \citep{emerging_threats} cung cấp hàng chục nghìn
luật công khai, trong đó các luật C2 thường neo vào \texttt{content}, \texttt{depth}
hoặc \texttt{offset} trong payload. Vì luật là tĩnh, một biến đổi làm dịch chuyển hoặc
phân mảnh các byte mục tiêu có thể khiến luật không còn khớp — đây là cơ sở kỹ thuật
cho các cơ chế phân mảnh mà chúng tôi dùng. \citet{silva2017ids} khảo sát né tránh IDS
bằng thao tác mạng dưới góc nhìn học máy.

\section{Học tăng cường và né tránh}

\MDP{} là khung hình thức chuẩn cho quyết định tuần tự \citep{puterman1994markov};
\PPO{} \citep{schulman2017proximal} là thuật toán policy-gradient ổn định, ít phân kỳ
off-policy, phù hợp với tập dữ liệu nhỏ. Thư viện Stable-Baselines3
\citep{stable_baselines3} và Gymnasium \citep{gymnasium} cung cấp hiện thực thực tế;
phiên bản \texttt{MaskablePPO} cho phép \emph{action masking} — chặn trước các hành
động không hợp lệ theo giao thức, giải quyết vấn đề khám phá trong không gian hành
động lớn.

\section{Khoảng trống mà luận văn hướng tới}

Phần lớn công trình trước dùng \emph{mô hình xấp xỉ} của IDS để chấm điểm, hoặc chạy
offline, hoặc không tính chi phí biến đổi. Chúng tôi lấp khoảng trống bằng cách:
(i) chấm bằng \emph{IDS thật} cho kết quả chính; (ii) luôn kèm \emph{control âm tính}
để chứng minh né tránh là thật; và (iii) báo cáo cả \emph{kết quả phủ định} khi RL
không tạo thêm giá trị — điều mà nhiều báo cáo né tránh không nêu.

% ===========================================================================
\chapter{Thiết kế hệ thống \& Mô hình mối đe dọa}
\label{ch:design}

\section{Mô hình mối đe dọa}

\subsection{Kẻ tấn công (red team)}
Kẻ tấn công là người vận hành botnet, nắm được nội dung gói tin C2 dạng plaintext
(PDU lệnh C2). Nó biết \emph{protocol framing} của luồng mình (HTTP hay chuỗi bản ghi
nhị phân) nhưng \textbf{không} biết nội dung bộ luật của defender. Nó chỉ được sửa
\textbf{payload} (dữ liệu TCP/UDP), không được sửa metadata (IP, cổng).

\subsection{Defender (blue team)}
Defender là \Snort{} 2.9.20 với bộ luật ET Open C2 (21\,374 luật), chạy tĩnh, không
phản ứng thích ứng. Mọi phán quyết (alert / không alert) được đọc trực tiếp từ nhị
phân Snort.

\subsection{Ràng buộc}
\begin{itemize}
  \item \textbf{Ngân sách hành động:} mỗi luồng chọn \emph{một} cơ chế biến đổi từ 14
        cơ chế; một số cơ chế chỉ áp dụng được cho một framing nhất định.
  \item \textbf{Giữ ngữ nghĩa:} phép biến đổi phải để lệnh C2 tới endpoint nguyên vẹn
        (định nghĩa ở mục~\ref{sec:semantics}).
  \item \textbf{Mục tiêu:} tối đa số luồng né được, tối thiểu chi phí biến đổi.
\end{itemize}

\section{Môi trường \texttt{HiddenDefenderEnv}}

Môi trường được hiện thực trong \texttt{ai\_agent/hidden\_defender\_env.py}. Một
\emph{episode} = một luồng = một quyết định.

\subsection{Quan sát (observation)}
Agent chỉ thấy 8 đặc trưng thống kê (black-box), \textbf{không} thấy bộ luật:
\begin{center}
\begin{tabular}{ll}
\toprule
Chỉ số & Đặc trưng \\
\midrule
0 & \texttt{is\_udp} (luồng UDP hay không) \\
1 & \texttt{n\_pkts\_norm} (số gói, chuẩn hoá) \\
2 & \texttt{tot\_bytes\_norm} (tổng byte, chuẩn hoá) \\
3 & \texttt{n\_payload\_pkts\_norm} (số gói có payload) \\
4 & \texttt{payload\_mean\_norm} (trung bình kích thước payload) \\
5 & \texttt{payload\_std\_norm} (độ lệch chuẩn payload) \\
6 & \texttt{payload\_max\_norm} (payload lớn nhất) \\
7 & \texttt{payload\_min\_norm} (payload nhỏ nhất) \\
\bottomrule
\end{tabular}
\end{center}

\subsection{Không gian hành động}
14 cơ chế (\texttt{ACTION\_NAMES}), chỉ giữ lại những cơ chế đã được \emph{đo} là
hoặc né được hợp lệ, hoặc là đối chứng âm hữu ích:

\begin{center}
\begin{tabular}{llr}
\toprule
Chỉ số & Cơ chế & Chi phí (\texttt{ACTION\_COSTS}) \\
\midrule
0 & \texttt{split8}  & 0{,}60 \\
1 & \texttt{split16} & 0{,}50 \\
2 & \texttt{split24} & 0{,}40 \\
3 & \texttt{prepend4}  & 0{,}20 \\
4 & \texttt{prepend8}  & 0{,}20 \\
5 & \texttt{prepend12} & 0{,}20 \\
6 & \texttt{pad16}   & 0{,}10 \\
7 & \texttt{reorder} & 0{,}35 \\
8 & \texttt{ttl}     & 0{,}10 \\
9 & \texttt{corrupt8} & 0{,}40 \\
10 & \texttt{overlap8} & 0{,}35 \\
11 & \texttt{noop}     & 0{,}00 \\
12 & \texttt{http\_header\_pad} & 0{,}05 \\
13 & \texttt{length\_wrapper}   & 0{,}25 \\
\bottomrule
\end{tabular}
\end{center}

\texttt{split8/16/24}, \texttt{prepend4/8/12} là các cơ chế phân mảnh/chèn đã được đo;
\texttt{pad16}, \texttt{reorder}, \texttt{ttl} được đo là 0/80 né (giữ làm đối chứng
âm có thông tin); \texttt{corrupt8} cố ý phá ngữ nghĩa (để hình phạt $-10$ quan sát
được); \texttt{noop} là hành động rỗng; \texttt{http\_header\_pad} và
\texttt{length\_wrapper} là hai cơ chế \emph{nhận biết giao thức} (gọi là A và B).

\subsection{Action masking (lọc hành động khả thi)}
Hàm \texttt{mech\_applicable()} chặn các cơ chế mà framing của luồng không thể mang:
\texttt{http\_header\_pad} chỉ áp dụng cho framing \texttt{http};
\texttt{length\_wrapper} chỉ cho framing \texttt{length2}. Các cơ chế không áp dụng
được \textbf{bị loại khỏi mẫu số}, để mọi tỷ lệ phần trăm đều tính trên các cơ chế
\emph{có thể} tác động lên luồng đó. Điểm quan trọng: mask là một \emph{bộ lọc tính
khả thi}, không phải một phép biến đổi — \texttt{MaskablePPO} đặt logit của hành động
bị mask về $-\infty$ trước softmax, nên hành động không khả thi không bao giờ được
lấy mẫu.

\subsection{Phần thưởng (reward)}
\begin{center}
\begin{tabular}{lr}
\toprule
Tình huống & Phần thưởng \\
\midrule
Né được \textbf{và} ngữ nghĩa còn nguyên & $+10$ (\texttt{REWARD\_EVADE}) \\
Né được \textbf{nhưng} phá ngữ nghĩa & $-10$ (\texttt{REWARD\_BREAK\_SEMANTICS}) \\
Vẫn bị báo động & $-1$ (\texttt{REWARD\_ALERTED}) \\
\bottomrule
\end{tabular}
\end{center}
Cộng thêm một phần thưởng nhỏ \texttt{PROTOCOL\_MATCH\_BONUS} $= 0{,}15$ \emph{chỉ khi}
một biến đổi đã né hợp lệ bằng \texttt{http\_header\_pad} trên luồng HTTP (nên nó
không thể ``chế ra'' né tránh hay che một cơ chế đang thất bại), và trừ
\texttt{ACTION\_COSTS} theo cơ chế. Phần thưởng
dương quan sát được trong kết quả chính là $9{,}6$ (với \texttt{split24}) và $9{,}4$
(với \texttt{split8}) — đúng bằng $10$ trừ chi phí tương ứng ($0{,}4$ và $0{,}6$).

\subsection{Định nghĩa ``giữ ngữ nghĩa''}
\label{sec:semantics}
Hàm \texttt{semantics\_intact()} yêu cầu lệnh C2 phải tới endpoint nguyên vẹn. Vì
không mô hình hoá được ngữ pháp của mọi giao thức, chúng tôi dùng mức chặt:
\begin{itemize}
  \item Nếu số gói thay đổi (phân mảnh), so sánh \emph{luồng byte được ghép lại}:
        lệnh parse được từ luồng gốc phải bằng lệnh parse được từ luồng biến đổi.
  \item Nếu số gói giữ nguyên, với \emph{mọi} gói mà payload gốc sinh ra một lệnh,
        lệnh đó phải bằng nhau trước và sau.
\end{itemize}
Đây là mức ``byte-equality'' chặt. Một mức lỏng hơn (chỉ kiểm tra chuỗi con) từng
khiến mọi cơ chế đều được tính là né (xem lỗi số~3, Chương~\ref{ch:discussion}).

% ===========================================================================
\chapter{Phương pháp}
\label{ch:method}

\section{Dữ liệu và phân chia tập}

\subsection{Nguồn dữ liệu}
Corpus lấy từ các capture malware thật của Stratosphere IPS
\citep{stratosphere_ctu13}. Hai capture được dùng cho tập huấn luyện và test:
\begin{itemize}
  \item \texttt{botnet-capture-20110811-neris} (Stratosphere MCFP)
  \item \texttt{botnet-capture-20110819-bot} (CTU-13)
\end{itemize}
Cả hai capture đều \emph{có mặt trong cả train và test}, để một phương pháp không thể
``thắng'' bằng cách nhận diện ``đây là capture nào''.

\subsection{Quy tắc chọn và phân chia}
Chỉ giữ lại các luồng mà Snort thật \emph{thật sự} phát hiện (để có cái mà né).
Phân chia cố định bởi \texttt{corpus\_config.py}:
\begin{center}
\begin{tabular}{lr}
\toprule
Tham số & Giá trị \\
\midrule
\texttt{n\_train} & 1\,153 \\
\texttt{n\_test}  & 40 \\
\texttt{SEED}     & 42 \\
\texttt{OBS\_DIM} & 8 \\
\bottomrule
\end{tabular}
\end{center}
Hàm \texttt{split\_flows()} dùng \texttt{random.Random(42)} để shuffle rồi lấy 40 luồng
đầu làm test — phân chia \emph{xác định} (deterministic) để tái lập được. Test cố ý
giữ ở 40 luồng để mọi eval/control trước đó vẫn so sánh được.

\section{Huấn luyện}

\subsection{Thuật toán}
Dùng \texttt{MaskablePPO} \citep{schulman2017proximal,stable_baselines3} với policy
\texttt{MlpPolicy}. Vì mỗi luồng là một quyết định độc lập nhưng PPO cần trajectory,
một wrapper mỏng (\texttt{BatchEnv}) gộp $n$ luồng liên tiếp thành \emph{một} episode
dài, để rollout là $n$ bước quyết định thật thay vì $n$ episode độ dài 1.

\begin{center}
\begin{tabular}{ll}
\toprule
Siêu tham số & Giá trị \\
\midrule
\texttt{n\_steps}   & $= n\_flows$ (mặc định 64) \\
\texttt{batch\_size}& $= n\_flows$ \\
\texttt{n\_epochs}  & 10 \\
\texttt{learning\_rate} & $3\times10^{-4}$ \\
\texttt{ent\_coef}  & 0{,}0151 \\
\texttt{gamma}      & 0{,}0 (mỗi luồng độc lập) \\
\texttt{seed}       & 42 \\
Ngưỡng dừng (\texttt{TARGET}) & 0{,}90 \\
\bottomrule
\end{tabular}
\end{center}
\texttt{gamma}=0 vì mỗi bước là một quyết định độc lập, không có phần thưởng tương lai.

\subsection{Hai bộ chấm điểm}
Đây là điểm \emph{quan trọng} về tính trung thực:
\begin{itemize}
  \item \textbf{Surrogate nhanh} (\texttt{RealRulesReplica}): parse bộ luật ET Open thật
        và tái hiện các sid thật sự kích hoạt; dùng để \emph{huấn luyện} cho nhanh
        ($\sim$0{,}33\,ms/luồng). Nó \textbf{bỏ sót 20{,}6\%} luật thật (các luật dùng
        \texttt{pcre}, \texttt{byte\_test}, \texttt{byte\_extract}, \texttt{byte\_jump}).
  \item \textbf{Snort thật} (\texttt{--real-snort}): chấm mọi hành động bằng nhị phân
        thật ($\sim$1{,}3\,s/luồng); dùng cho \emph{kết quả chính}.
\end{itemize}
Việc tách hai bộ chấm là có chủ đích: con số $99{,}2\%$ khi huấn luyện bằng surrogate
\emph{không phải} cùng đại lượng với con số $90{,}6\%$ khi huấn luyện bằng Snort thật.

\section{Quy trình đo lường}

\subsection{Một cặp (flow, cơ chế) cho mỗi pcap}
Gộp 32 cặp vào một pcap từng khiến Snort phát 4 alert trong khi chấm riêng lẻ tìm được
12 phát hiện — lưu lượng trộn lẫn làm giảm alert trên toàn capture. Vì vậy ma trận
cơ chế chạy với \texttt{batch\_size}=1: 224 lần gọi Snort, $\sim$80\,s.

\subsection{Control âm tính (bắt buộc)}
Mỗi lần đo đều kèm một control: các luồng \emph{chưa} biến đổi phải bị Snort báo động.
Nếu control ngừng báo động, mọi con số khác bị bỏ. Control này là thứ đã bắt được hai
lỗi tệ nhất của dự án.

\subsection{Trần oracle và các baseline}
Để con số ``PPO đạt 100\%'' có ý nghĩa, cần một trần:
\begin{center}
\begin{tabular}{lll}
\toprule
Phương pháp & Bản chất & Vai trò \\
\midrule
\texttt{random\_sweep}   & thứ tự cơ chế ngẫu nhiên, thử tới khi hợp lệ & sàn \\
\texttt{greedy}          & dùng lại cơ chế hiệu quả trên luồng gần nhất & heuristic mạnh \\
\texttt{supervised\_tree}& cây quyết định trên nhãn oracle & \textbf{trần} \\
\texttt{ppo}/\texttt{maskable\_ppo} & policy đã huấn luyện & đối tượng cần đánh giá \\
\bottomrule
\end{tabular}
\end{center}

\subsection{Verifier tự động}
\texttt{controls/verify\_report\_claims.py} tái suy diễn mọi con số trong báo cáo từ
file JSON gốc và \emph{thất bại} nếu có sai lệch. Verifier kiểm tra các \emph{bất biến}
(\eg{} \texttt{queries == epochs * n\_flows}) chứ không ghim hằng số, vì dừng sớm làm
hằng số thay đổi mỗi lần chạy. Bộ test (\texttt{tests/}) gồm 69 test đều pass.

% ===========================================================================
\chapter{Kết quả \& Đánh giá}
\label{ch:results}

Mọi con số trong chương này được đọc từ file JSON trong
\texttt{snort\_validation/reports/}. Không có con số nào được ước lượng.

\section{Kết quả chính (Snort thật)}

Từ \texttt{ppo\_masked\_test.json} — \texttt{eval\_masked\_ppo\_test.py --real-snort},
trên 40 luồng held-out:

\begin{center}
\begin{tabular}{lr}
\toprule
Chỉ số & Giá trị \\
\midrule
method & \texttt{maskable\_ppo} \\
scorer & \texttt{real\_snort} \\
split & test \\
$n\_flows$ & 40 \\
\textbf{alerted} & \textbf{0} \\
\textbf{valid\_evasion} & \textbf{40 / 40 (100{,}0\%)} \\
\texttt{masked\_actions\_used} & 0 \\
Control âm tính (\texttt{control\_noop\_real\_snort.json}) & \textbf{40 / 40 bị báo động} \\
\bottomrule
\end{tabular}
\end{center}

Policy chỉ dùng đúng hai cơ chế:
\begin{center}
\begin{tabular}{lr}
\toprule
Cơ chế & Số luồng \\
\midrule
\texttt{split8}  & 24 (toàn bộ luồng UDP/neris) \\
\texttt{split24} & 16 (toàn bộ luồng HTTP/TCP) \\
\bottomrule
\end{tabular}
\end{center}
Mọi luồng đều có \texttt{semantics\_ok = true}; phần thưởng là 9{,}4 (split8) và 9{,}6
(split24). Việc phân bố framing/action là hoàn hảo (24 neris$\to$split8, 16 http$\to$split24),
đúng như kỳ vọng về masking nhận biết giao thức.

\subsection{Control âm tính}
\texttt{control\_noop\_real\_snort.json}: 40 luồng \emph{không} biến đổi, tất cả
\texttt{alert = true}, \texttt{evaded = false}. Tỷ lệ né của control là \textbf{0\%}.
Đây là bằng chứng định lượng rằng 40/40 của agent là thật, không phải do Snort không
chạy.

\section{Số liệu huấn luyện}

\begin{center}
\begin{tabular}{llrrrrl}
\toprule
Lần chạy & Bộ chấm & flows & né được & queries & epochs & file \\
\midrule
Surrogate & \texttt{RealRulesReplica} & 128 & 127 (99{,}2\%) & 896 & 7  & \texttt{ppo\_hidden\_defender.json} \\
Snort thật & nhị phân thật & 64 & 58 (90{,}6\%) & 640 & 10 & \texttt{ppo\_realsnort\_train.json} \\
\bottomrule
\end{tabular}
\end{center}
\textbf{Cả hai lần đều dừng ở ngưỡng $\ge 90\%$, không phải ở hội tụ.} Không lần nào
là trần thật.

\section{Ma trận cơ chế (Snort thật, 8 luồng HTTP)}

Từ \texttt{hidden\_defender\_matrix\_test.json} — 8 luồng HTTP, 112 ô, chấm một cặp
mỗi pcap:

\begin{center}
\begin{tabular}{lrr}
\toprule
Cơ chế & Né được & Né được \textbf{và} ngữ nghĩa nguyên \\
\midrule
\texttt{split8}          & 8/8 & \textbf{8/8} \\
\texttt{split16}         & 8/8 & \textbf{8/8} \\
\texttt{split24}         & 8/8 & \textbf{8/8} \\
\texttt{http\_header\_pad} & 8/8 & \textbf{8/8} \\
\texttt{prepend4}        & 8/8 & 0/8 \\
\texttt{prepend8}        & 8/8 & 0/8 \\
\texttt{prepend12}       & 8/8 & 0/8 \\
\texttt{overlap8}        & 8/8 & 0/8 \\
\texttt{length\_wrapper} & 8/8 & 0/8 \\
\texttt{corrupt8}        & 1/8 & 0/8 \\
các cơ chế còn lại      & 0/8 & 0/8 \\
\bottomrule
\end{tabular}
\end{center}

Đây là bảng quan trọng nhất vì cột phải là nơi công việc trung thực thể hiện:
\textbf{sáu cơ chế đạt 8/8 né nhưng 0/8 né hợp lệ}. Chúng không đánh bại IDS — chúng
phá kênh C2, và một hàm thưởng ngây thơ sẽ cho chúng điểm tuyệt đối.
\texttt{length\_wrapper} còn là code chết: 0/8 hợp lệ \emph{và} 0/16 áp dụng được trên
luồng neris.

\section{RL có thật sự cần thiết không? (phát hiện làm phức tạp câu chuyện)}

Từ \texttt{hidden\_defender\_real\_snort\_test.json} — 8 luồng HTTP, Snort thật:

\begin{center}
\begin{tabular}{lrr}
\toprule
Phương pháp & Queries & Né được \\
\midrule
\texttt{random\_sweep}    & 25 & 8/8 \\
\texttt{greedy}           & 13 & 8/8 \\
\texttt{supervised\_tree} & 8  & 8/8 \\
\textbf{ppo}              & 8  & 8/8 \\
\bottomrule
\end{tabular}
\end{center}

Trên tập huấn luyện 64 luồng, so sánh còn rõ hơn: \texttt{greedy} né 58/64 với
\textbf{149} queries, trong khi PPO cũng né 58/64 nhưng tốn \textbf{704} queries —
gấp \textbf{4{,}7 lần}. Và trên tập HTTP, bốn cơ chế \texttt{split8}/\texttt{split16}/
\texttt{split24}/\texttt{http\_header\_pad} đều giải được 8/8, nên hành động tối ưu
gần như \emph{hằng số theo framing}, không phải theo luồng.

\textbf{Cách đặt vấn đề trung thực cho báo cáo:} đây \emph{không phải} ``RL đánh bại
các baseline''. Mà là: ``môi trường là thật, phép đo là vững, và dưới \emph{defender
này} bài toán tìm kiếm quá dễ để RL tạo thêm giá trị — hành động thắng gần như cố
định.'' Đây là một \emph{kết quả phủ định về cách đóng khung bài toán}, không phải về
agent, và nó bảo vệ được.

\section{Scale-up (Snort thật, neris-20110811, 8 vòng)}

Từ \texttt{final\_scale\_up\_summary.json}:

\begin{center}
\begin{tabular}{rrrr}
\toprule
Số luồng & Né được & Tỷ lệ né & mean corrupt \\
\midrule
24  & 24  & 100{,}0\% & 0{,}00 \\
100 & 94  & 94{,}0\%  & 8{,}69 \\
300 & 287 & 95{,}7\%  & 7{,}71 \\
323 & 69  & \textbf{21{,}4\%}  & 0{,}00 \\
\bottomrule
\end{tabular}
\end{center}

Kết quả \textbf{không đơn điệu}: từ 24 lên 300 luồng vẫn cao, rồi sụp ở 323 luồng.
Chúng tôi nêu thẳng hiện tượng này thay vì làm mượt nó (xem Chương~\ref{ch:discussion}).

\section{Kiểm tra chéo capture (cross-capture)}

Từ \texttt{cross\_capture\_summary.json} (cost $=0{,}6$), sinh lại từ các báo cáo
per-capture:

\begin{center}
\begin{tabular}{lrrrr}
\toprule
Capture & Baseline bị bắt & Policy né & Corrupt-all né & Kết luận \\
\midrule
botnet-capture-20110810-neris & 23/24 & 4  & 4  & hoà \\
botnet-capture-20110811-neris & 24/24 & 24 & 24 & hoà \\
capture-win13                 & 24/24 & 0  & 0  & hoà \\
\bottomrule
\end{tabular}
\end{center}

\textbf{Trung bình 38{,}89\%, độ lệch chuẩn 43{,}74\%}, min 0\%, max 100\%.
\texttt{vs\_control\_counts = \{tie: 3\}}.

Cột ``Corrupt-all né'' là điểm mấu chốt: trên cả ba capture, policy né đúng bằng số
luồng mà hành động ``làm hỏng mọi gói có thể'' cũng né được, với cùng mức biến đổi
trung bình. \textbf{Policy hoà control 3/3} — nó không học được gì mà hành động tầm
thường ``corrupt-everything'' chưa đạt được. Vì vậy con số 100\% trên capture 0811 là
\emph{tính chất của control đủ mạnh}, không phải bằng chứng cho RL agent.
\texttt{capture-win13} ở 0\% là tín hiệu trung thực rằng policy không tổng quát hoá
qua capture.

\section{Kiểm định surrogate so với Snort thật}

Từ \texttt{surrogate\_sweep\_14mech.json} — 14 cơ chế $\times$ 16 luồng, 224 lần gọi
Snort thật, 2\,359\,s:

\begin{center}
\begin{tabular}{lr}
\toprule
Chỉ số & Giá trị \\
\midrule
Đồng thuận tổng thể & 161/224 (71{,}9\%) \\
Đồng thuận trên ô áp dụng được & 103/116 (88{,}8\%) \\
\textbf{surrogate false-safe} (báo né khi Snort báo động) & \textbf{0} \\
\bottomrule
\end{tabular}
\end{center}
\texttt{surrogate\_vs\_real\_snort.json}: 8/8 đồng thuận, 0 false-safe. Con số ``0
false-safe'' là con số quan trọng nhất: một false-safe nghĩa là surrogate nói dối
agent rằng nó đã né — đúng chế độ thất bại có thể vô hiệu hoá mọi thí nghiệm chấm bằng
surrogate. Nó \emph{không} xảy ra ở đây.

\section{Sweep chi phí corrupt}
Từ \texttt{sweep\_corrupt\_cost.json}: né 100\% từ cost 0{,}2 đến 2{,}0, rồi 0\% ở
cost 4{,}0; mean corrupt bị ghim ở 2{,}88 qua mọi mức cost. Tham số cost không định
hình policy mà chỉ bật/tắt nó.

% ===========================================================================
\chapter{Thảo luận}
\label{ch:discussion}

\section{Vì sao \texttt{MaskablePPO} hợp ở đây}

\begin{itemize}
  \item \textbf{Action masking giải quyết khám phá.} Với framing, cơ chế khả thi bị
        lọc trước, nên agent chỉ cân nhắc các biến đổi hợp lệ theo giao thức.
  \item \textbf{\PPO{} ổn định trên corpus nhỏ.} \PPO{} tránh phân kỳ off-policy, phù
        hợp khi dữ liệu ít (1\,153 luồng).
  \item \textbf{Quan sát black-box 8 chiều.} Agent không thấy bộ luật, chỉ thấy thống
        kê luồng — đúng như một kẻ tấn công thật.
\end{itemize}

\section{Chi phí và ý nghĩa}

Mỗi biến đổi có chi phí (\texttt{ACTION\_COSTS}): phân mảnh tốn băng thông/CPU, chèn
header HTTP rẻ và tự nhiên. Trong phiên bản bandit đã nghỉ hưu, chi phí được hình thức
hoá qua \texttt{break\_even\_cost} $\approx 3{,}2$ gói: với $N$ gói có thể sửa trong
một luồng, ``làm hỏng mọi gói'' chỉ hoà vốn khi chi phí vượt \texttt{EVASION\_BONUS}/$N$.
Với môi trường hiện tại, chi phí là các hằng số cố định theo cơ chế; điều đáng chú ý
là sweep chi phí \emph{không} định hình policy (mean corrupt ghim ở 2{,}88) — đây là
một dấu hiệu overfit cần nêu thẳng.

\section{Vì sao các con số nên được tin: những lỗi đã bắt được}

Mỗi lỗi sau từng tạo ra kết quả sai trông rất hợp lý, và bị bắt bởi control hoặc
kiểm tra chéo — không phải bởi đọc code:

\begin{center}
\begin{tabular}{cll}
\toprule
\# & Lỗi & Cách phát hiện \\
\midrule
1 & Trích xuất luồng một chiều (mất gói ngược) & baseline phát hiện 0/24 (bất khả) \\
2 & Viết lại gói tin ngây thơ (sửa cả phản hồi server) & vẫn 0/24 sau khi sửa lỗi 1 \\
3 & \texttt{semantics\_intact} chấp nhận chứa chuỗi con & mọi cơ chế đạt 100\% \\
4 & Chấm điểm theo lô làm triệt alert & pcap 32 cặp cho 4 alert, chấm riêng 12 \\
5 & Luật ET throttle \texttt{by\_src} & địa chỉ query tái dùng alert một lần rồi im \\
6 & Cửa sổ đọc alert ngắn hơn nhịp flush của Snort & \texttt{--resident} báo baseline 0/24 \\
\bottomrule
\end{tabular}
\end{center}

\textbf{Lỗi số~3 là lỗi cần trình bày cẩn thận nhất.} Nó là khác biệt giữa ``agent né
được IDS'' và ``agent né được IDS \emph{trong khi vẫn giữ malware hoạt động}''. Kết quả
100\% trước đó đã bị thu hồi vì nó. Mức hiện tại là byte-equality.

\section{Hạn chế}

\begin{itemize}
  \item \textbf{Chỉ một seed, một cặp capture.} Kết quả chính là một seed trên tập test
        40 luồng, tất cả từ hai capture cùng họ. Chưa có trần hội tụ đa-seed.
  \item \textbf{Defender tĩnh, đơn luật.} Chỉ bộ luật ET Open tĩnh; chưa thử IDS thích
        ứng hay phát hiện bất thường.
  \item \textbf{Surrogate bỏ sót 20{,}6\% luật.} Trên corpus này không ảnh hưởng (chỉ 3
        sid kích hoạt mỗi capture) nhưng sẽ ảnh hưởng trên lưu lượng khác. Đồng thuận
        \texttt{split8} chỉ 5/16 (31{,}2\%) — surrogate tệ nhất đúng nơi policy hoạt động.
  \item \textbf{Không tổng quát hoá chéo capture.} win13 ở 0\%. Đây là ngõ cụt của
        \emph{không gian hành động}: các luồng win13 gần như không có payload (0/97 gói
        có \texttt{Raw}), nên không byte nào sửa được; Snort vẫn alert 24/24 vì luật
        establishment/dsize không cần payload. Kết luận đúng là phụ thuộc \emph{trường
        luật}, không phải phụ thuộc họ botnet.
  \item \textbf{Scale không đơn điệu.} Sụp ở 323 luồng chưa được giải thích; cần chẩn
        đoán hoặc nói thẳng ``chúng tôi chưa biết''.
  \item \textbf{Cả hai cơ chế thắng đều là phân mảnh} (\texttt{split8}, \texttt{split24})
        — policy chọn cách né rẻ nhất còn giữ ngữ nghĩa, không phải một kho cơ chế đa
        dạng.
\end{itemize}

\section{Các bất nhất cần giải quyết}

Tài liệu \texttt{docs/REPORT\_WRITING\_GUIDE.md} và
\texttt{docs/SIC\_AI\_FRAMEWORK\_VALIDATION.md} còn mang một số con số \emph{đã lỗi
thời}: ví dụ mean evasion 73{,}3\%, control 71/71, train 1\,113 luồng, không gian hành
động ``65\,536'', phần thưởng $+100/-2/-50$. Các con số này \textbf{không} khớp với
file JSON gốc. Luận văn này dùng \emph{file JSON gốc} làm nguồn chân lý (xem Phụ lục~E).
Một số tài liệu khác (\texttt{all\_models\_comparison.json},
\texttt{mechanism\_matrix.json}) thuộc thế hệ hệ thống cũ và không so sánh được với kết
quả hiện tại.

\section{Hướng phát triển}

\begin{itemize}
  \item \textbf{Defender dị thể/thích ứng:} cấp cho mỗi luồng một luật C2 với
        \texttt{depth}/\texttt{offset} khác nhau, để offset phân mảnh thắng thay đổi
        thật sự theo luồng — khi đó RL mới có việc để làm.
  \item \textbf{Đa seed và kiểm tra chéo họ botnet thật} (không chỉ cùng nhị phân).
  \item \textbf{Retrain surrogate theo kiểu DAgger} (gán nhãn chính biến đổi của agent
        mới bằng Snort thật rồi huấn luyện lại) để đóng khoảng cách surrogate--thực tế.
  \item \textbf{Mở rộng không gian hành động} sang metadata/TTL/phân mảnh để không bị
        ngõ cụt trên các luồng không payload như win13.
\end{itemize}

% ===========================================================================
\chapter{Kết luận}
\label{ch:conclusion}

Luận văn cho thấy một agent \texttt{MaskablePPO} có thể biến đổi gói tin C2 thật để né
\Snort{} 2.9.20 (ET Open) một cách đáng tin cậy: \textbf{40/40 luồng held-out né hợp lệ
dưới Snort thật}, trong khi control âm tính bị báo động 40/40. Mọi hành động đều được
chấm bằng nhị phân thật, và một verifier tự động 81 kiểm tra bảo đảm không con số nào
trôi dạt khỏi nguồn.

Nhưng luận văn cũng kết luận trung thực rằng, \emph{dưới defender tĩnh này}, RL chưa
tạo thêm giá trị: bài toán tìm kiếm gần như suy biến (một hành động gần như cố định
mỗi framing), nên PPO không vượt greedy; và trên kiểm tra chéo capture (trung bình
38{,}89\%, std 43{,}74\%), policy \emph{hoà} control ``corrupt-everything'' trên cả ba
capture — con số né chéo capture không phải thành tựu của RL. Một capture đạt 0\% vì lý
do không gian hành động, không phải vì họ botnet.

Giá trị của công trình nằm ở \emph{sự trung thực của phép đo}: control âm tính, trần
oracle, một cặp mỗi pcap, verifier tự động, và việc báo cáo thẳng cả kết quả phủ định.
Con đường phía trước rõ ràng: làm defender dị thể/thích ứng để bài toán thật sự cần
đến RL, rồi kiểm tra đa seed và chéo họ botnet.

% ===========================================================================
\chapter{Phụ lục}
\label{ch:appendix}

\section{A. Thống kê corpus}
\begin{center}
\begin{tabular}{lr}
\toprule
Chỉ số & Giá trị \\
\midrule
Tổng luồng train & 1\,153 \\
Tổng luồng test  & 40 \\
\texttt{obs\_dim} & 8 \\
Capture train & \texttt{20110811-neris}, \texttt{20110819-bot} \\
Capture test  & \texttt{20110811-neris}, \texttt{20110819-bot} \\
Framing trong tập test & 24 \texttt{neris} (UDP) + 16 \texttt{http} (TCP) \\
\bottomrule
\end{tabular}
\end{center}

\section{B. Chi tiết hiện thực}
\textbf{Snort:} 2.9.20, bộ luật ET Open C2 (21\,374 luật), lệnh
\texttt{snort -c <conf> -r <pcap> -l <dir> -q -A fast}; phán quyết đọc từ file fast-alert.
Chấm một cặp (flow, cơ chế) mỗi pcap. \textbf{Thời gian:} eval chính \texttt{--real-snort}
$\sim$172\,s, control $\sim$166\,s, ma trận cơ chế $\sim$11{,}2\,s/ô (112 ô
$\approx$20 phút). \textbf{Hyperparameters:} xem Chương~\ref{ch:method}.
\textbf{Verifier:} 81/81 kiểm tra pass. \textbf{Test:} 69 test pass.

\section{C. Bảng kết quả per-flow (40 luồng test)}
Bảng~\ref{tab:perflow} liệt kê toàn bộ 40 luồng held-out: framing, cơ chế được chọn,
alert, ngữ nghĩa, phần thưởng. Tất cả đều \texttt{alert=false}, \texttt{evaded=true},
\texttt{semantics\_ok=true}.

\begin{table}[h]
\centering
\small
\caption{Kết quả per-flow trên tập test held-out (Snort thật).}
\label{tab:perflow}
\begin{tabular}{rlllrr}
\toprule
\# & Capture & Framing & Action & Reward & Sem. \\
\midrule
0  & 20110811-neris & http  & split24 & 9{,}6 & \checkmark \\
1  & 20110819-bot  & neris & split8  & 9{,}4 & \checkmark \\
2  & 20110811-neris & http  & split24 & 9{,}6 & \checkmark \\
3  & 20110819-bot  & neris & split8  & 9{,}4 & \checkmark \\
4  & 20110819-bot  & neris & split8  & 9{,}4 & \checkmark \\
5  & 20110819-bot  & neris & split8  & 9{,}4 & \checkmark \\
6  & 20110819-bot  & neris & split8  & 9{,}4 & \checkmark \\
7  & 20110819-bot  & neris & split8  & 9{,}4 & \checkmark \\
8  & 20110819-bot  & neris & split8  & 9{,}4 & \checkmark \\
9  & 20110819-bot  & neris & split8  & 9{,}4 & \checkmark \\
10 & 20110819-bot  & neris & split8  & 9{,}4 & \checkmark \\
11 & 20110811-neris & http  & split24 & 9{,}6 & \checkmark \\
12 & 20110819-bot  & neris & split8  & 9{,}4 & \checkmark \\
13 & 20110819-bot  & neris & split8  & 9{,}4 & \checkmark \\
14 & 20110811-neris & http  & split24 & 9{,}6 & \checkmark \\
15 & 20110819-bot  & neris & split8  & 9{,}4 & \checkmark \\
16 & 20110819-bot  & neris & split8  & 9{,}4 & \checkmark \\
17 & 20110811-neris & http  & split24 & 9{,}6 & \checkmark \\
18 & 20110819-bot  & neris & split8  & 9{,}4 & \checkmark \\
19 & 20110811-neris & http  & split24 & 9{,}6 & \checkmark \\
20 & 20110819-bot  & neris & split8  & 9{,}4 & \checkmark \\
21 & 20110811-neris & http  & split24 & 9{,}6 & \checkmark \\
22 & 20110811-neris & http  & split24 & 9{,}6 & \checkmark \\
23 & 20110811-neris & http  & split24 & 9{,}6 & \checkmark \\
24 & 20110811-neris & http  & split24 & 9{,}6 & \checkmark \\
25 & 20110811-neris & http  & split24 & 9{,}6 & \checkmark \\
26 & 20110819-bot  & neris & split8  & 9{,}4 & \checkmark \\
27 & 20110819-bot  & neris & split8  & 9{,}4 & \checkmark \\
28 & 20110819-bot  & neris & split8  & 9{,}4 & \checkmark \\
29 & 20110819-bot  & neris & split8  & 9{,}4 & \checkmark \\
30 & 20110811-neris & http  & split24 & 9{,}6 & \checkmark \\
31 & 20110811-neris & http  & split24 & 9{,}6 & \checkmark \\
32 & 20110819-bot  & neris & split8  & 9{,}4 & \checkmark \\
33 & 20110811-neris & http  & split24 & 9{,}6 & \checkmark \\
34 & 20110819-bot  & neris & split8  & 9{,}4 & \checkmark \\
35 & 20110811-neris & http  & split24 & 9{,}6 & \checkmark \\
36 & 20110819-bot  & neris & split8  & 9{,}4 & \checkmark \\
37 & 20110819-bot  & neris & split8  & 9{,}4 & \checkmark \\
38 & 20110811-neris & http  & split24 & 9{,}6 & \checkmark \\
39 & 20110819-bot  & neris & split8  & 9{,}4 & \checkmark \\
\bottomrule
\end{tabular}
\end{table}

\section{D. Control âm tính và kiểm tra ngữ nghĩa}
Control \texttt{control\_noop\_real\_snort.json}: 40 luồng, tất cả bị báo động
(\texttt{alerted=40}, \texttt{n=40}), tỷ lệ né 0\%. Mức ngữ nghĩa dùng trong môi trường
là byte-equality / so khớp lệnh sau khi ghép luồng (mục~\ref{sec:semantics}). Chỉ khi
\texttt{evaded AND semantics\_ok} thì một luồng mới được tính là né hợp lệ
(\texttt{QueryCounter.valid()}).

\section{E. Ghi chú chỉnh sửa các con số lỗi thời}
Một số tài liệu nội bộ từng mang con số \emph{đã bị vượt qua}. Bảng dưới đối chiếu
con số trong tài liệu với con số trong file JSON gốc (nguồn chân lý):

\begin{center}
\begin{tabular}{lll}
\toprule
Chỉ số & Giá trị cũ (tài liệu) & Giá trị đúng (JSON gốc) \\
\midrule
Mean evasion chéo capture & 73{,}3\% / 62{,}5\% & \textbf{38{,}89\%} \\
Std chéo capture          & 10{,}8\% / 44{,}49\% & \textbf{43{,}74\%} \\
Control                   & 71/71 bị bắt & \textbf{40/40 bị bắt} \\
Tập test né được          & 29/40 & \textbf{40/40} \\
Luồng train               & 1\,113 & \textbf{1\,153} \\
Không gian hành động      & 65\,536 & \textbf{14 cơ chế} \\
Phần thưởng               & $+100/-2/-50$ & \textbf{$+10/-10/-1$} \\
\bottomrule
\end{tabular}
\end{center}
Nguyên nhân gốc của dữ liệu tổng hợp lỗi thời: script tổng hợp cũ lọc input theo chuỗi
tên file \texttt{"summary"} và bỏ sót khi cập nhật, để lại một file cũ. Hiện tại script
chọn input theo \emph{schema} và verifier khẳng định file tổng hợp khớp với mọi báo cáo
per-capture, nên một file tổng hợp lỗi thời sẽ làm verifier \emph{thất bại} thay vì
được công bố.

% ===========================================================================
\bibliographystyle{plainnat}
\bibliography{thesis/references}

\end{document}
